%% ============================================================
%% 片段：s02a-definition   对应原文 PDF 第 5--13 页
%% 本文件由 gw/convert.py 从 arXiv 源码（hep-th/0612073v2）自动生成：
%%   * 公式/编号/引用/图 已经锁定为原文内容，请勿改动任何数学内容与编号；
%%   * 只需把英文正文译为中文，并在每节开头加“本节重点”、关键处加译注。
%% ============================================================


\section{单值性与面算符}

%% ==== 原文 §2：PDF 第 5 页（印刷第 4 页）====
\label{sec:2}

\begin{keytakeaways}
  \begin{itemize}[leftmargin=1.4em,itemsep=2pt,topsep=2pt]
    \item \textbf{物理图像}：面算符由这样的要求定义：在四维 $\mathcal N=4$ 超杨–米尔斯理论中，场沿余维二子流形 $D$ 取奇异的剖面；在垂直于 $D$ 的二维法平面上，该剖面是希钦方程的一个奇异解。这是 't Hooft 算符（经由余维三奇点）在余维二上的类比。
    \item \textbf{数学工具}：法平面上的希钦方程 $F-\phi\wedge\phi=0$、$d_A\phi=0$、$d_A\star\phi=0$；解空间的模空间 $\MH$ 是超 Kähler 流形（即带分歧的希格斯丛的模空间），其 Kähler 类与复结构分别由参数 $(\alpha,\beta,\gamma)$ 控制（超 Kähler 商的构造见第~\ref{sec:3}~节）。
    \item \textbf{关键结论}：在复结构 $I$ 中面算符给出希格斯丛（$\beta+i\gamma$ 为复模、$\alpha$ 为 Kähler 模），在 $J$ 中给出平坦 $G_{\mathbb{C}}$-丛的单值性（$\gamma+i\alpha$ 为复模、$\beta$ 为 Kähler 模），在 $K$ 中给出 $\alpha+i\beta$ 复模、$\gamma$ 为 Kähler 模；见表~\ref{tab:1}（原文 Table 1）。
    \item \textbf{离散数据与参数}：$\alpha$ 只定义到余特征格平移与外尔变换（二者合起来是仿射外尔群的作用）；$\alpha$ 的量子化条件为 $f(\alpha)\,D\cap D\in\Z$（式~\eqref{eq:2.9}、\eqref{eq:2.10}；$U(1)$ 情形即 $\alpha D\cap D\in\Z$）。量子化后还出现二维 $\theta$ 角 $\eta\in{}^L\TT$，它在 $S$-对偶下与 $\alpha$ 互换（见第~\ref{sec:2.4}~节）。
  \end{itemize}
\end{keytakeaways}


\subsection{面算符的定义}
 %% ==== 原文 §2.1：PDF 第 5 页（印刷第 4 页）====
\label{sec:2.1}

本文的基本技术，是研究 $\mathcal N=4$ 超杨–米尔斯理论（super Yang–Mills theory）如何能沿时空中的一个余维二子流形被修改。因此，我们考虑四维流形 $M$ 上的 $\mathcal N=4$ 超杨–米尔斯理论，但沿一个二维子流形 $D$ 作修改，使四维场沿 $D$ 出现奇点。于是，这一构造是通常经由余维三奇点构造 't Hooft 算符之做法在面算符（surface operator）上的类比。一般性背景可参见 \cite{ref38} 或 \cite{ref8} 第 6.2 节。

\begin{yizhu}[物理图像补充]
这个类比只到“用奇点定义算符”为止，两者的参数性质恰好相反：'t Hooft 算符是余维三的，环绕它的小 $S^2$ 上的磁通量必须量子化，故磁荷是离散的，取值于格 $\Lambda_{\rm cochar}$；面算符是余维二的，法平面里只剩绕 $D$ 的圆周，$\alpha$ 只由绕 $D$ 的和乐 $\exp(-2\pi\alpha)$ 决定，因而是连续的，属于环面 $\TT$。面算符真正的离散数据是 $D$ 上 $\TT$-丛的拓扑类 $\mathfrak m\in\Lambda_{\rm cochar}$，它到第~\ref{sec:2.3}~节才出场，并与 $\theta$ 角 $\eta$ 配对；也正因 $\alpha$ 连续，$S$-对偶才能把它与 $\eta$ 互换（第~\ref{sec:2.4}~节）。
\end{yizhu}

我们的重点将是 $\mathcal N=4$ 超杨–米尔斯理论的 GL 扭曲（GL-twisted）版本，它是几何朗兰兹（geometric Langlands）纲领的规范场论进路的基础。规范群（gauge group）是紧李群 $G$，我们一般假定它是单群。最重要的玻色场，一为规范场（gauge field）$A$，它是 $G$-丛 $E\to M$ 上的联络（connection）；一为取值于 ${\rm ad}(E)$ 的一形式场 $\phi$。我们的规范场论约定与 \cite{ref8} 相同。特别地，$A$ 与 $\phi$ 取值于 $G$ 的实李代数（因而在 $G$ 的酉表示中由反厄米矩阵表示），协变导数为 $D=d_A=d+A$，和乐（holonomy）为 $P\exp\left(-\int A\right)$。

沿 $D$ 出现奇点的场，正是 $A$ 与 $\phi$ 到 $D$ 的法丛（normal bundle）上的限制。局部地，可以把四维流形模型化为 $M=D\times D'$，其中 $D'$ 是法丛的纤维，而奇点位于 $D'$ 中的一点，例如原点。

GL 扭曲 $\mathcal N=4$ 超杨–米尔斯理论的超对称方程依赖于一个参数 $t$，如 \cite{ref8} 所解释。但约化到二维后，这一参数消失，剩下的就是希钦方程（Hitchin's equations）：

% 式 (2.1)
\setcounter{equation}{0}
\begin{equation}\label{eq:2.1}
  \begin{aligned}F-\phi\wedge \phi& = 0 \\
                     d_A\phi & = 0 \\
                     d_A\star \phi & = 0.\end{aligned}
\end{equation}
 （此处 $\star$ 为霍奇星（Hodge star）算符。）
\begin{yizhu}[易错点]
希钦方程的几种等价写法在此后各节反复出现，须与本节约定对齐：本节取物理文献约定 $F-\phi\wedge\phi=0$、$d_A\phi=0$、$d_A\star\phi=0$，其中 $\phi$ 是实一形式、$A$ 反厄米。在法平面的复坐标 $z=x^1+ix^2$ 下写 $\phi=\varphi+\bar\varphi$，则 $d_A\phi=0$ 与 $d_A\star\phi=0$ 合起来等价于 Dolbeault 形式 $\bar\partial_A\varphi=0$（即 $\varphi$ 全纯），而 $F-\phi\wedge\phi=0$ 给出 $F$ 的 $(1,1)$ 部分与 $[\varphi,\bar\varphi]$ 成正比的矩映射型方程。若改用 $F+\phi\wedge\phi=0$ 或把 $\phi$ 换成 $i\phi$ 的约定，则式~\eqref{eq:2.11} 中的组合 $\beta+i\gamma$ 与表~\ref{tab:1} 的对应关系随之改变，引用本文结果前务必先核对 $\phi$ 的归一与符号。
\end{yizhu}

因此，我们将通过描述希钦方程的解中可能出现的一种孤立奇点来定义面算符。

事实上，本文只需要最简单的一类奇点。先取 $D'=\mathbb{R}^2$，其欧氏坐标为 $x^1,x^2$，使得 $\star (dx^1)=dx^2$、$\star(dx^2)=-dx^1$。再引入极坐标 $x^1+ix^2=re^{i\theta}$。记 $\TT$ 为 $G$ 的极大环面（maximal torus），并记 $\mathfrak g$、$\mathfrak t$ 分别为 $G$ 与 $\TT$ 的李代数。为描述在原点具有孤立奇点的希钦方程解，取元素 $\alpha,\beta,\gamma\in \mathfrak t$，并令

% 式 (2.2)
\setcounter{equation}{1}
\begin{equation}\label{eq:2.2}
  \begin{aligned}A &
= \alpha \,d\theta \\
           \phi & = \beta\,{dr\over r}-\gamma\,d\theta.\end{aligned}
\end{equation}

由于 $\alpha$、$\beta$、$\gamma$ 彼此交换（$\mathfrak t$ 是交换的），且一形式 $d\theta$ 与 $dr/r$ 既闭又余闭，希钦方程成立。我们完全按照通常定义 't Hooft 算符的方式来定义面算符（例如参见 \cite{ref38} 或 \cite{ref8} 第 6.2 节）：要求在 $r=0$ 附近，$A$ 与 $\phi$ 的行为如式~\eqref{eq:2.2}。在整体情形下，对二维子流形 $D\subset M$，我们要求：在 $D$ 的每一点，法平面内的场（在某个规范下）都形如式~\eqref{eq:2.2}，且 $\alpha$、$\beta$、$\gamma$ 取指定的值。

\begin{yizhu}[含糊点拨]
式~\eqref{eq:2.2} 为何自动满足式~\eqref{eq:2.1}，可以直接验算：由 $\star(dx^1)=dx^2$、$\star(dx^2)=-dx^1$ 得 $\star\,dr=r\,d\theta$、$\star\,(dr/r)=d\theta$、$\star\,d\theta=-dr/r$，故 $d\theta$ 与 $dr/r$ 在 $\mathbb{R}^2\backslash\{0\}$ 上既闭又余闭。又因 $\alpha,\beta,\gamma$ 两两交换，$A\wedge\phi-\phi\wedge A=0$ 且 $[\phi,\phi]=0$，而 $F=dA=0$（$r>0$ 处），三条方程便逐条成立。换成复坐标 $z=re^{i\theta}$ 时 $dz/z=dr/r+i\,d\theta$，$dr/r$ 与 $d\theta$ 恰好互为霍奇对偶；于是 $\phi=\beta\,dr/r-\gamma\,d\theta$ 分解为 ${1\over2}(\beta+i\gamma)\,dz/z+{1\over2}(\beta-i\gamma)\,d\bar z/\bar z$，其 $(1,0)$ 部正是式~\eqref{eq:2.11} 中的留数 $\beta+i\gamma$ 乘以 $dz/z$。注意 $d\theta$ 前的负号是必需的：若写成 $+\gamma\,d\theta$，$(1,0)$ 部就变成 ${1\over2}(\beta-i\gamma)\,dz/z$ 了。
\end{yizhu}

显然，可以用 $G$ 的外尔群（Weyl group）${\mathcal W}$ 作用于三元组 $(\alpha,\beta,\gamma)$，而不本质地改变理论。因此，我们所定义的面算符只依赖于 $(\alpha,\beta,\gamma)$ 模一个外尔变换。但还有额外的自由度。若 $u\in \mathfrak t$ 满足 $\exp(2\pi u)=1$，则由 $\TT$-值函数

% 式 (2.3)
\setcounter{equation}{2}
\begin{equation}\label{eq:2.3}
  (r,\theta)\to \exp(\theta u)
\end{equation}

 给出的规范变换使 $\alpha$ 平移 $u$。模去这一变换后，$\alpha$ 的惟一不变量是联络 $A$ 绕常 $r$ 圆周的单值性（monodromy）；该单值性为 $\exp(-2\pi \alpha)$。于是，三元组 $(\alpha,\beta,\gamma)$ 取值于 $\TT\times {\mathfrak t}\times {\mathfrak t}$，模去 ${\mathcal W}$ 的作用。
\begin{yizhu}[易错点]
单值性的符号约定是全篇最容易出错之处，务必与 \cite{ref8} 对齐：本节取 $A$ 反厄米、和乐为 $P\exp(-\int A)$，故 $A=\alpha\,d\theta$ 绕原点一周给出 $\exp(-2\pi\alpha)$，而不是 $\exp(2\pi\alpha)$，也不是 $\exp(-2\pi i\alpha)$——$\alpha$ 是 $d\theta$ 的系数，指数上不带 $i$。由此，第~\ref{sec:2.2}~节中 $\CA=(\alpha-i\gamma)d\theta$ 给出 $U=\exp(-2\pi(\alpha-i\gamma))$，而 $\tilde\CA=A+i\star\phi$ 给出 $\exp(-2\pi(\alpha+i\beta))$；这些 $2\pi$、$i$ 与 $\pm$ 都由同一和乐约定固定，并正好与表~\ref{tab:1} 的三个复模 $(\beta+i\gamma,\gamma+i\alpha,\alpha+i\beta)$ 一一对应。若改用正厄米约定（$A\to iA$、和乐取 $P\exp(i\int A)$），则全部单值性变为原来的复共轭，$\alpha,\beta,\gamma$ 的实性条件也需重设。
\end{yizhu}

\begin{yizhu}[易错点]
三元组取值于 $\TT\times\mathfrak t\times\mathfrak t$ 而不是 $\TT^3$，根源在规范变换对三者的作用不同。取式~\eqref{eq:2.3} 的 $g=\exp(\theta u)$，则 $g^{-1}dg=u\,d\theta$ 只把 $A$ 变为 $(\alpha+u)\,d\theta$，而 $\phi\to g\phi g^{-1}=\phi$ 因 $u$ 与 $\beta,\gamma$ 交换而完全不变。于是 $\beta,\gamma$ 没有格平移，只有外尔群作用，参数空间为 $\mathfrak t/{\mathcal W}$；$\alpha$ 还须再商 $\Lambda_{\rm cochar}$，即 $\TT/{\mathcal W}$（完整商是仿射外尔群 ${\mathcal W}_{\rm aff}={\mathcal W}\ltimes\Lambda_{\rm cochar}$）。这也说明 $\beta$ 与 $\gamma$ 分工不同：$\beta$ 是 $dr/r$ 的系数，$\gamma$ 是 $d\theta$ 的系数，前者给出极点留数的实部、后者给出虚部，二者在式~\eqref{eq:2.11} 中合成为同一个复留数 $\beta+i\gamma$。
\end{yizhu}

我们将经常、但并非总是对 $\alpha$ 使用加法记号。这相当于把 $\TT$ 视为对某个格 $\Lambda$ 的商 ${\mathfrak t}/\Lambda$。为确定这个格，注意若 $\TT={\mathfrak t}/\Lambda$，则 $\Lambda=\pi_1(\TT)={\rm Hom}(U(1),\TT)$。我们称之为 $G$ 的余特征格（cocharacter lattice），记作 $\Lambda_{\rm cochar}$。（更多信息见附录 A。）事实上，${\rm Hom}(U(1),\TT)$ 参数化 $\TT$ 值的磁荷（magnetic charge），或等价地——由基本的 GNO 对偶 \cite{ref39}——参数化对偶群 $^L\neg G$ 的电荷。$\Lambda_{\rm cochar}$ 是余权格（coweight lattice）$\Lambda_{\rm cowt}$ 的子格。二者之商是 $G$ 的中心 ${\mathcal Z}(G)$：

% 式 (2.4)
\setcounter{equation}{3}
\begin{equation}\label{eq:2.4}
  {\mathcal Z}(G)=\Lambda_{\rm cowt}/\Lambda_{\rm cochar}.
\end{equation}

\begin{yizhu}[含糊点拨]
${\rm Hom}(U(1),\TT)$ 之所以能参数化磁荷，是因为它同构于 $\pi_1(\TT)$：$\TT$-丛在 $S^2$ 上的拓扑类由绕赤道的单值性给出，而 $\pi_1(\TT)=\Lambda_{\rm cochar}$。对 $G=SU(2)$，$\Lambda_{\rm cochar}=\mathbb{Z}$ 正是 't Hooft–波利亚科夫单极子的磁荷；GNO 对偶把这个格解释为对偶群 $^L\neg G$ 的电荷格，于是式~\eqref{eq:2.4} 的商也有了读法：$G$ 单连通时 $\Lambda_{\rm cochar}$ 是余根格、$\Lambda_{\rm cowt}$ 是余权格，商即中心 ${\mathcal Z}(G)$（$SU(N)$ 时为 $\mathbb{Z}_N$）。$G$ 不单连通时两个格都要按 $G$ 重算，原文仍断言式~\eqref{eq:2.4}。
\end{yizhu}


\bigskip\noindent{\it 丛的扩张}


先来尝试计算奇异联络 $A=\alpha\,d\theta$ 在原点处的曲率。由于 $d(d\theta)=2\pi
\delta_D$（$\delta_D$ 是支撑在 $D'$ 原点、庞加莱对偶于 $D$ 的二形式 delta 函数），我们似乎得到

% 式 (2.5)
\setcounter{equation}{4}
\begin{equation}\label{eq:2.5}
  F=2\pi\alpha \delta_D.
\end{equation}
 然而这不可能是自然的公式，因为我们有把 $\alpha$ 平移一个格矢的自由。曲率的定义出了什么问题？我们把 $A$ 作为 $G$-丛 $E$ 上的联络引入，但由于沿 $D$ 的奇点，该丛只在 $M$ 中 $D$ 的补集上自然有定义。可以选取 $E$ 到 $D$ 上的一个扩张，但没有自然的扩张。$E$ 在 $D$ 上不同的可能扩张，对应于把 $\alpha$ 从其自然取值所在的 $\TT={\mathfrak t}/\Lambda_{\rm cochar}$ 提升到 $\mathfrak t$ 的不同方式。一旦选定了扩张，原点处曲率的计算才有意义，结果即为式~\eqref{eq:2.5}。

使 $\alpha$ 平移一个格矢的规范变换 $(r,\theta)\to \exp(\theta u)$，由于其在原点处的奇点，把 $E$ 在 $D$ 上的一个扩张映为另一个扩张。尽管没有把 $E$ 作为 $G$-丛扩张到 $D$ 上的自然方式，我们仍可以这样做。在 $D$ 附近，$E$ 的结构群自然地约化到与 $A$ 和 $\phi$ 的奇异部分交换的子群 $\TT$。奇异规范变换 $\theta\to \exp(\theta u)$ 在 $\TT$ 上平凡地作用；因此，尽管 $D$ 上没有自然的 $G$-丛，$D$ 上却有自然的 $\TT$-丛。我们将假定 $A$ 在 $D$ 上的限制就是这个 $\TT$-丛上的联络，并假定曲率 $F$ 限制到 $D$ 时同样取值于 $\mathfrak t$。

假设规范群 $G$ 不是单连通的；例如它可以是伴随型的群。此时规范变换 $\exp(\theta u)$ 可能不会提升为 $G$ 的单连通覆盖 $\bar G$ 中的单值规范变换；在 $\theta\to\theta+2\pi$ 下，它会被乘以 $\bar G$ 的中心 ${\mathcal Z}(\bar G)$ 中的一个元素 $y$。当出现这种情形时，该规范变换会改变丛 $E$ 的拓扑，其方式是平移特征类（characteristic class）$\xi\in H^2(M,\pi_1(G))$——后者度量把 $E$ 提升为 $\bar G$-丛的障碍。若对 ${\mathcal Z}(\bar G)$ 采用加法记号，则平移为

% 式 (2.6)
\setcounter{equation}{5}
\begin{equation}\label{eq:2.6}
  \xi\to \xi+y[D],
\end{equation}
 其中 $[D]$ 是庞加莱对偶于 $D$ 的类。于是，具有不同 $\xi$ 值、以及适当相关的 $\alpha$ 值的规范场论彼此等价。

这一构造还有一个变体如下。假设规范群实际上就是 $\bar G$，或是其它使由 $\exp(\theta u)$ 给出的规范变换不再单值的群。这个规范变换作为 $\mathcal N=4$ 超杨–米尔斯理论（其中所有场都在伴随表示中）的对称性，局部地仍然有意义。若整体拓扑使得一个在 $D$ 附近局部地形如 $\exp(\theta u)$ 的规范变换可以整体延拓到 $M$ 上，则 $\mathcal N=4$ 超杨–米尔斯理论在 $\alpha\to\alpha+u$ 下不变。这里 $u$ 可以是 $\Lambda_{\rm cowt}$ 中满足 $y[D]=0$ 的任意元素。更一般地，若 $D$ 是互不相交的分支 $D_i$ 的并，而在各 $D_i$ 附近作由 $\exp(\theta_i u_i)$ 给出的规范变换、且 $y_i=\exp(2\pi u_i)$，则条件是

% 式 (2.7)
\setcounter{equation}{6}
\begin{equation}\label{eq:2.7}
  \sum_iy_i[D_i]=0.
\end{equation}



\bigskip\noindent{\it 非平凡法丛}

在引出这一构造时，我们从乘积 $M=D\times D'$ 这一特例出发。然而更一般地，对任意嵌入的二维流形 $D\subset M$，我们考虑在每个法平面内具有式~\eqref{eq:2.2} 型奇点的规范场。为简单起见，本文只考虑 $M$ 与 $D$ 均可定向的情形。$D$ 可以有非平凡法丛，从而有非零的自交数（self-intersection number）$D\cap D$，后者可刻画为\footnote{关于这一点以及随后立刻要作出的若干断言，参见 \cite{ref40} 中关于 Thom 类的描述。}
% 式 (2.8)
\setcounter{equation}{7}
\begin{equation}\label{eq:2.8}
  D\cap D
=\int_M\delta_D\wedge \delta_D.
\end{equation}
 当自交数非零时，$\alpha$ 不可能整体地取任意值。先对 $G=U(1)$ 说明这一点。在 $D$ 附近具有奇点 $A=\alpha d\theta$ 的联络 $A$ 具有性质 $\int_D F/2\pi = \alpha D\cap D~{\rm mod}~\Z$。由于积分后的第一陈类 $\int_D F/2\pi$ 总是整数，故总有

% 式 (2.9)
\setcounter{equation}{8}
\begin{equation}\label{eq:2.9}
  \alpha D\cap D\in \Z.
\end{equation}
 对任意 $G$，其推广不过是：同样的陈述在 $\TT$ 的任意 $U(1)$ 子群中成立。因此，若 $\alpha\to f(\alpha)$ 是 $\mathfrak t$ 上任一在格 $\Lambda_{\rm cochar}$ 上取整数值的实值线性函数，则
% 式 (2.10)
\setcounter{equation}{9}
\begin{equation}\label{eq:2.10}
  f(\alpha) D\cap D\in\Z.
\end{equation}
\begin{yizhu}[易错点]
自交数非零时 $\alpha$ 的量子化条件 (2.9)--(2.10) 有两个容易漏掉的前提：其一，此时已选定 $E$ 到 $D$ 的一个扩张，$\int_D c_1(E)=\langle c_1(E),[D]\rangle$ 才是 $c_1$ 与 $[D]$ 的整配对；其二，(2.10) 在 $\alpha\to\alpha+u$（$u\in\Lambda_{\rm cochar}$）下保持不变，因为 $f(u)\in\Z$ 而 $D\cap D\in\Z$——这正是第~\ref{sec:2.1}~节末尾所说的“只有与 (2.10) 相容的平移才能整体定义”的代数根源，也是第~\ref{sec:2.4}~节中 $\alpha$ 与二维 $\theta$ 角 $\eta$ 可以互换的技术前提。另外，当 $D\cap D\ne0$ 时，二维 $\theta$ 角的定义还须在作用量中加入正比于 $D\cap D\,\Tr\alpha\eta$ 的 $c$-数项，配分函数不再是 $\TT\times{}^L\TT$ 上的函数而是复线丛的截面；几何朗兰兹的通常情形取 $D\cap D=0$，故这些效应可以略去。
\end{yizhu}

同样正确的是：若 $D\cap D\not= 0$，则在法平面中由式~\eqref{eq:2.3} 定义的扭曲规范变换（twisted gauge transformation）并不总能沿 $D$ 整体地定义。只有那些按与式~\eqref{eq:2.10} 相容的方式平移 $\alpha$ 的规范变换，才能真正整体地定义。

沿满足 $D\cap D\not=0$ 的曲面的奇点在四维流形理论 \cite{ref11} 中很重要，但在本文的应用中不那么重要，因为几何朗兰兹纲领按其通常形式处理的正是 $D\cap D=0$ 的情形。

\subsection{参数的几何解释}

%% ==== 原文 §2.2：PDF 第 10 页（印刷第 9 页）====
\label{sec:2.2}

我们通过要求场在 $D$ 附近的行为如式~\eqref{eq:2.2}，定义了支撑在二维流形 $D\subset M$ 上的面算符。一般而言，在量子情形下，希钦方程、乃至理论的二阶经典场方程都不必被满足；面算符的定义只要求它们在 $D$ 附近被满足。然而，为更好地理解参数 $\alpha,\beta,\gamma$ 在经典几何中的含义，我们考虑这样的情形：确实要在黎曼面（Riemann surface）$C$（它对应于上述分析中的 $D'$）上求解希钦方程，并在某点 $p\in C$ 附近带一个上述类型的孤立奇点。（这里只解释为引出第~\ref{sec:2.4}~节的对偶猜想所需的经典几何事实；第~\ref{sec:3}~节将以远为详尽的方式重新讨论经典几何。）

具有式~\eqref{eq:2.2} 所述类型奇点的希钦方程解，已在 \cite{ref22,ref23,ref24} 中分析过。与希钦方程的光滑解一样，这类解的模空间是超 Kähler（hyper-Kähler）流形 $\EUM_H$，我们称之为带分歧的希格斯丛（ramified Higgs bundle）的模空间，也称具有抛物结构（parabolic structure）的希格斯丛的模空间。（若想显式指明规范群、黎曼面或奇点的性质与位置，我们把它记作 $\EUM_H(G)$、$\EUM_H(C)$、$\EUM_H(\alpha,\beta,\gamma;p)$ 等。）

由于超 Kähler 结构，希钦方程的解可以从不同角度看待。从一个通常称为 $I$ 的复结构（complex structure）的角度看，黎曼面 $C$ 上希钦方程的一个解描述一个希格斯丛（Higgs bundle），即一对 $(E,\varphi)$，其中 $E$ 是全纯（holomorphic）$G$-丛，$\varphi$ 是 $K_C\otimes
{\rm ad}(E)$ 的全纯截面。（$K_C$ 记 $C$ 的典范丛（canonical bundle）。）希格斯丛的构造如下，从希钦方程的解 $(A,\phi)$ 出发。把规范协变外微分 $d_A=d+A$ 的 $(0,1)$ 部分解释为 $\bar\partial_A$ 算符，它赋予丛 $E$ 全纯结构；并把 $\varphi$ 定义为 $\phi$ 的 $(1,0)$ 部分，也就是把 $\phi$ 分解为 $\varphi+\bar\varphi$，其中 $\varphi$ 为 $(1,0)$ 型、$\bar\varphi$ 为 $(0,1)$ 型。于是希钦方程蕴含 $\varphi$ 全纯，且 $(E,\varphi)$ 是希格斯丛。
\begin{yizhu}[原文疑误]
原文源码第 800 行作 \texttt{of type (1,0}，右括号落在数学环境之外（源码里数学环境的起止符位置错置）。PDF 第 10 页上这一错置几乎不可见（括号仍大致排在同一位置），但同一页及后文其余各处均写作 $(1,0)$；译文按 $(1,0)$ 排出，不改其数学含义。
\end{yizhu}

在当前情形下，令 $z=x^1+ix^2$，由式~\eqref{eq:2.2} 可得

% 式 (2.11)
\setcounter{equation}{10}
\begin{equation}\label{eq:2.11}
  \varphi = {1\over 2}(\beta+i\gamma){dz\over z}.
\end{equation}

于是，从复结构 $I$ 的角度看，面算符在希格斯场中引入一个留数（residue）为 $(1/2)(\beta+i\gamma)$ 的极点\footnote{该陈述在 $\beta+i\gamma$ 是李代数 ${\mathfrak g}_\C$ 的“正则”元素时成立（与之交换的 ${\mathfrak g}_\C$ 的子代数恰为 ${\mathfrak t}_\C$）。其他情形见第~\ref{sec:3.3}~节。}。$\varphi$ 的特征多项式（characteristic polynomial）在复结构 $I$ 中全纯地变化，特别地，$\varphi$ 中极点的共轭类（conjugacy class）也如此。式~\eqref{eq:2.11} 表明，该极点的共轭类全纯地依赖于 $\beta+i\gamma$，且与 $\alpha$ 无关。这是一项更一般陈述的一部分：复结构 $I$ 随 $\beta+i\gamma$ 全纯地变化，且与 $\alpha$ 无关。另一方面，相应的 Kähler 形式 $\omega_I$ 在复结构 $I$ 中为 $(1,1)$ 型，而其上同调类与 $\beta$、$\gamma$ 无关，且是 $\alpha$ 的线性函数。这些陈述以及下面若干类似的陈述，都来自把模空间构造为超 Kähler 商（hyper-Kähler quotient）\cite{ref23,ref24}，我们将在第~\ref{sec:3.2}~节解释。因此，若从复结构 $I$ 的视角看 $\MH$，则 $\beta+i\gamma$ 是复参数（$I$ 全纯地依赖于它），而 $\alpha$ 是 Kähler 参数。
\begin{yizhu}[含糊点拨]
作者在此以“这些陈述来自超 Kähler 商构造 [23,24]”一句带过，略去了为何三个参数恰好按表~\ref{tab:1} 分配。其实质是：把 $\MH$ 看作超 Kähler 商，其三个矩映射的水平（level）中，实部决定辛形式（从而决定 Kähler 类），复部决定复结构；$\alpha$ 只进入实水平，故只改变 Kähler 类，而 $\beta+i\gamma$ 只进入复水平，故只改变复结构。三者在 $I,J,K$ 之间的 $SO(3)$ 旋转下按三重态变换，于是得到表~\ref{tab:1} 的循环对称；这也正是复结构 $I$ 全纯地依赖 $\beta+i\gamma$、却与 $\alpha$ 无关的原因。
\end{yizhu}

\begin{yizhu}[物理图像补充]
把这段话译成复几何的语言：复结构 $I$ 下的 $\MH$ 就是 $C$ 上带分歧的希格斯丛的模空间，分歧点 $p$ 处的抛物结构由 $\alpha$ 决定——$\alpha$ 在 $p$ 处定出一个博雷尔（或抛物）子代数与一组抛物权重，$\beta+i\gamma$ 则定出 $\varphi$ 在 $p$ 处极点的留数共轭类（式~\eqref{eq:2.11}）；$G=U(N)$ 时这组数据就是本节末所说 \cite{ref22} 表中的 $(\alpha,2b,2c)$。一般 $G$ 的版本见第~\ref{sec:3.4}~节：带奇点的平坦丛与抛物丛一一对应，$\alpha$ 决定抛物子群，$\alpha$ 的格平移对应赫克修改。因此“$\alpha$ 决定稳定性”与“$\alpha$ 给出抛物权”是同一件事的两种说法。
\end{yizhu}

事实上，尽管本文不需要这些细节，$\alpha$ 在纯代数几何中也有自然的含义（无需提及 Kähler 度量），但这个含义有些微妙。复结构 $I$ 中的 $\MH$ 参数化这样的对 $(E,\varphi)$：$\varphi$ 具有单极点（simple pole），其共轭类由 $\beta$ 与 $\gamma$ 决定（如式~\eqref{eq:2.11}），并且 $(E,\varphi)$ 是“稳定”的。相应的稳定性（stability）概念 \cite{ref22} 依赖于 $\alpha$。

\begin{yizhu}[易错点]
“单极点”（simple pole）这三个字不能省：式~\eqref{eq:2.11} 的 $dz/z$ 正是一阶极点；若换成高阶极点（野分歧，wild ramification），留数与稳定性的讨论都要改写，本文只处理这种最简单的情形。另需注意本节有两个不同的“正则”条件：脚注所说的是 $\beta+i\gamma$ 在 $\mathfrak g_\C$ 中正则（中心化子恰为 $\mathfrak t_\C$），保证式~\eqref{eq:2.11} 的留数共轭类有良定义；第~\ref{sec:2.3}~节开头所说的则是三元组 $(\alpha,\beta,\gamma)$ 在 $G$ 中正则（稳定化子恰为 $\TT$），保证沿 $D$ 的结构群确实约化到 $\TT$。两者失效的非正则情形都在第~\ref{sec:3.3}~节讨论。
\end{yizhu}

在复结构 $J$ 中，自然的复变量是 $G_{\mathbb{C}}$ 值的联络 $\CA=A+i\phi$；按希钦方程它是平坦的。$\CA=(\alpha-i\gamma)d\theta$ 绕 $p$ 处奇点的单值性为

% 式 (2.12)
\setcounter{equation}{11}
\begin{equation}\label{eq:2.12}
  U=\exp(-2\pi(\alpha-i\gamma)).
\end{equation}
 它全纯地依赖于 $\gamma+i\alpha$，且与 $\beta$ 无关。确实，在复结构 $J$ 中，$\gamma+i\alpha$ 是复参数，而 $\beta$ 是 Kähler 参数。

在复结构 $J$ 中，$\MH$ 参数化 $C$ 上这样的平坦 $G_{\mathbb{C}}$-丛：绕 $p$ 点的单值性落在包含 $U$ 的共轭类\footnote{如上一脚注，该陈述在 $U$ 为正则元素时成立；更一般情形的讨论推迟到第~\ref{sec:3.3}~节。}之中。与复结构 $I$ 中的 $\alpha$ 一样，$\beta$ 在复结构 $J$ 中可以用纯全纯的语言解释（无需提及 Kähler 度量），但这种解释有些难以捉摸（在本文中只会起很小的作用）。据 \cite{ref22}，在复结构 $J$ 中，$\beta$ 决定了 $C\backslash p$（即挖去点 $p$ 后的 $C$）上以 $\CA$ 为联络的平坦丛的一个单值性不变的加权滤过（weighted filtration）的权。

最后，在复结构 $K=IJ$ 中，自然的复变量是 $G_{\mathbb{C}}$ 值的联络 $\tilde\CA=A+i\star
\phi$。由希钦方程它同样是平坦的。它绕 $p$ 处奇点的单值性为 $\exp(-2\pi(\alpha+i\beta))$。在复结构 $K$ 中，$\alpha+i\beta$ 是复结构参数，而 $\gamma$ 是 Kähler 参数。$\gamma$ 在复结构 $K$ 中的解释，与 $\beta$ 在复结构 $J$ 中的解释完全类似。
\begin{table}[htbp]
  \centering
  \begin{tabular}{c|c|c}
    \hline
    模型 & 复模 & K\"ahler 模 \\
    \hline
    $I$ & $\beta+i\gamma$ & $\alpha$ \\
    $J$ & $\gamma+i\alpha$ & $\beta$ \\
    $K$ & $\alpha+i\beta$ & $\gamma$ \\
    \hline
  \end{tabular}
  \caption{$\MH$ 在复结构 $I$、$J$、$K$ 中的复模与 K\"ahler 模：复结构 $I$ 全纯地依赖于 $\beta+i\gamma$，而相应的 K\"ahler 结构依赖于 $\alpha$。（原文 Table 1，PDF 第 12 页）}
  \label{tab:1}
\end{table}

\begin{yizhu}[易错点]
表~\ref{tab:1} 的三行须与正文逐条对应。$I$ 行取自式~\eqref{eq:2.11} 的留数 ${1\over2}(\beta+i\gamma)$，去掉因子 $1/2$ 得 $\beta+i\gamma$（正文随后说明“为得到循环对称性，我们调整了一个因子 2”）；$J$ 行取自式~\eqref{eq:2.12} 的 $\exp(-2\pi(\alpha-i\gamma))$，而表中写 $\gamma+i\alpha$，二者相差因子 $i$，因为 $\gamma+i\alpha=i(\alpha-i\gamma)$；$K$ 行取自单值性 $\exp(-2\pi(\alpha+i\beta))$，表中写 $\alpha+i\beta$。乘上非零复常数不改变“复结构是否全纯依赖于该参数”，故三种写法都合法，但循环对称 $(I,J,K)\leftrightarrow(\beta+i\gamma,\gamma+i\alpha,\alpha+i\beta)$ 只在这样取时才成立；若把 $J$ 行照抄成 $\alpha-i\gamma$，对称性即被破坏。
\end{yizhu}


这些陈述总结于表~\ref{tab:1}。表中可见，在同时置换 $I,J,K$ 与 $\alpha,\beta,\gamma$ 时存在循环对称性。（这一循环对称性是出现在一个密切相关情形中的 $SO(3)$ 对称性的一部分；见第~\ref{sec:3.8}~节。）若 $G=U(N)$，则 $(\alpha,\beta,\gamma)$ 在 \cite{ref22} 第 720 页的表中被称为 $(\alpha,2b,2c)$。为得到循环对称性，我们调整了一个因子 2。


\subsection{\texorpdfstring{$\theta$ 角}{θ 角}}

%% ==== 原文 §2.3：PDF 第 12 页（印刷第 11 页）====
\label{sec:2.3}

在量子情形下，除 $\alpha$、$\beta$、$\gamma$ 之外还存在一个额外参数。可以猜测它会如此，因为到此为止我们在每个复结构中只描述了一个实 Kähler 参数（列于表~\ref{tab:1} 的最后一列），而在超对称理论中，Kähler 参数通常是复化（complexified）的。

为说明这一点，暂假定三元组 $(\alpha, \beta,\gamma)$ 是正则的（regular），即保持该三元组不变的 $G$ 的子群恰为 $\TT$。要求每个法平面内都具有式~\eqref{eq:2.2} 的行为，特别意味着沿 $D$ 给出了丛 $E$ 的结构群从 $G$ 到 $\TT$ 的约化。因此，沿 $D$ 我们得到的是规范群为 $\TT$ 的阿贝尔规范场论（abelian gauge theory）。（我们在第~\ref{sec:2.1}~节末尾说明过：沿 $D$ 存在自然的 $\TT$-丛，尽管没有自然的 $G$-丛。）

\begin{yizhu}[易错点]
这里出现的是面算符的第二种描述，不要把它当成另一个对象：第~\ref{sec:2.1}~节用“四维场在法平面内的奇异配置”（式~\eqref{eq:2.2}）定义面算符，本节则说沿 $D$ 只剩下规范群为 $\TT$ 的二维规范场论。两者是同一对象的两个视角：前者是 $M$ 上的局部条件，后者是它诱导在 $D$ 上的有效自由度（$\TT$-丛的拓扑类 $\mathfrak m$ 与 $\theta$ 角 $\eta$）。还要注意“结构群从 $G$ 约化到 $\TT$”是定义的一部分（由三元组的正则性保证），不是额外假设；忽略这一点，就会把式~\eqref{eq:2.2} 的奇点与沿 $D$ 的二维理论误当成两种不同的构造。
\end{yizhu}

在二维阿贝尔规范场论中，“$\theta$ 角（theta angle）”起重要作用。例如取规范群为 $U(1)$。$U(1)$-丛 ${\mathcal L}\to D$ 在拓扑上由其次数（degree）$d=\int_D c_1({\mathcal L})$ 分类，其中 $c_1({\mathcal L})$ 是第一陈类（first Chern class）\footnote{为简单起见，假定 $D$ 是闭的。否则，在定义量子场论时需要给 $D$ 的边界（或无穷远）加上适当的边界条件；经过适当处理，仍可给出 $c_1({\mathcal L})$ 与 $\theta$ 的合适定义。}。$\theta$ 角通过一个相位 $\exp(i\theta d)$ 进入理论，该相位作为额外因子包含在路径积分中。这里 $\theta$ 取值于 $\mathbb{R}/2\pi\mathbb{Z}\cong U(1)$。

\begin{yizhu}[易错点]
本节有大小两个 $\theta$ 角，切勿混为一谈：一个是 $D$ 上的二维 $\theta$ 角 $\eta\in{}^L\neg\TT$（式~\eqref{eq:2.13}），有 $r$ 个实分量；另一个是四维规范场论的 $\theta$ 角，它与耦合合成 $\tau=\theta/2\pi+4\pi i/g^2$（第~\ref{sec:2.4}~节，该节取 ${\rm Im}\,\tau=4\pi/e^2$），故 $\theta\to\theta+2\pi$ 恰是 $T:\tau\to\tau+1$。二维一侧的 $\theta$ 只通过 $\exp(i\theta\mathfrak m)$ 进入配分函数（$\mathfrak m\in\Lambda_{\rm cochar}$ 为 $D$ 上 $\TT$-丛的类），这正是下文把 $\eta$ 写成同态 $\Lambda_{\rm cochar}\to U(1)$ 的理由；Witten 效应则把这个 $2\pi$ 与 $\alpha$ 的平移联系起来，见第~\ref{sec:2.5}~节式~\eqref{eq:2.50}。
\end{yizhu}

现在回到我们的问题，即非阿贝尔规范场论中的面算符。设 $G$ 的秩为 $r$。那么其极大环面 $\TT$ 同构于 $U(1)^r$，而规范群为 $G$ 的二维规范场论将有 $r$ 个 $\theta$ 角，取值于一个 $r$ 维环面。来看看这究竟是哪个环面。从拓扑观点看，二维流形 $D$ 上的 $\TT$-丛可以如下唯一地构造：从次数为 1 的 $U(1)$ 丛出发，再经由某个同态 $\rho:U(1)\to \TT$ 把它映入 $\TT$。因此，$D$ 上的 $\TT$-丛在拓扑上由取值于格 $\Lambda_{\rm cochar}={\rm Hom}(U(1),\TT)$ 的特征类（characteristic class）$\mathfrak{m}$ 分类。所以，$\TT$ 规范场论的 $\theta$ 角是一个同态 $\eta:\Lambda_{\rm cochar}\to U(1)$。

换言之，$\eta$ 取值于 ${\rm Hom}(\Lambda_{\rm cochar},U(1))$。我们断言

% 式 (2.13)
\setcounter{equation}{12}
\begin{equation}\label{eq:2.13}
  {\rm Hom}(\Lambda_{\rm cochar},U(1))={}^L\negthinspace \TT,
\end{equation}

 其中 $^L\neg \TT$ 是对偶群（dual group）$^L\neg G$ 的极大环面。事实上，由庞特里亚金对偶（Pontryagin duality）\footnote{庞特里亚金对偶断言：若 $W$ 是局部紧阿贝尔群，$V={\rm Hom}(W,U(1))$，则 $W={\rm Hom}(V,U(1))$。在我们的应用中，$W=\Lambda_{\rm cochar}={\rm Hom}(U(1),\TT)$，而 $V={}^L\neg \TT$。}，并利用 $\Lambda_{\rm cochar}={\rm Hom}(U(1),\TT)$，式~\eqref{eq:2.13} 等价于

% 式 (2.14)
\setcounter{equation}{13}
\begin{equation}\label{eq:2.14}
  {\rm Hom}(U(1),\TT)={\rm
Hom}(^L\neg \TT,U(1)).
\end{equation}
 而这是群与其对偶群之间关系的一个标准刻画。确实，左端分类 $G$ 的磁荷，右端分类 $^L\neg G$ 的电荷。两者的相等就是基本的 GNO 对偶 \cite{ref39}。

正如 $\TT={\mathfrak t}/\Lambda_{\rm cochar}$，我们有 $^L\neg
\TT={}^L{\mathfrak t}/\Lambda_{\rm char}$，其中 ${}^L\mathfrak t$ 是 ${}^L\neg \mathbb{T}$ 的李代数，而 $\Lambda_{\rm char}={\rm
Hom}(T,U(1))$ 是 $G$ 的电荷格，即特征格（character lattice）。${}^L\mathfrak t$ 与 $\mathfrak t$ 的对偶 $\mathfrak t^\vee$ 重合。我们经常对 $\eta$ 使用加法记号，把它视为 ${\mathfrak t}^\vee/\Lambda_{\rm char}$ 的元素。
\begin{yizhu}[易错点]
本节只说 $\eta$ 取值于 ${\rm Hom}(\Lambda_{\rm cochar},U(1))={}^L\TT$，但未强调它与 $\alpha$（取值于 $\TT={\mathfrak t}/\Lambda_{\rm cochar}$）的对称地位。用加法记号，$\alpha\in{\mathfrak t}/\Lambda_{\rm cochar}$ 与 $\eta\in{\mathfrak t}^\vee/\Lambda_{\rm char}$ 之间有自然配对 $\Tr\alpha\eta$，这正是第~\ref{sec:2.4}~节作用量中 $\Tr\alpha{\mathfrak m}$ 与 $\Tr\alpha\eta$ 各项的来源。此外，$\alpha$ 的规范等价不只有式~\eqref{eq:2.3} 的格平移，还有外尔群作用；两者合成半直积 ${\mathcal W}\ltimes\Lambda_{\rm cochar}$ 在 $\mathfrak t$ 上的作用（$G$ 单连通时 $\Lambda_{\rm cochar}$ 即余根格，此即仿射外尔群 ${\mathcal W}_{\rm aff}$），故 $\alpha$ 的物理参数空间是 $\mathfrak t/{\mathcal W}_{\rm aff}$，等价地是 $\TT/{\mathcal W}$。又因 ${\mathcal W}(G)={\mathcal W}({}^L G)$（见附录 A），同一外尔群也作用在 $\eta$ 上；第~\ref{sec:2.5}~节由 $T:\theta\to\theta+2\pi$ 导出的平移 $\eta\to\eta-\alpha$ 与这一外尔群作用对易，二者合起来才是作用在 $(\alpha,\eta)$ 上的完整对称性（当 $D\cap D\ne0$ 时还须补上与 $D\cap D$ 成正比的 $c$-数项）。
\end{yizhu}

\begin{yizhu}[含糊点拨]
作者说式~\eqref{eq:2.13} 与式~\eqref{eq:2.14} 由庞特里亚金对偶等价，具体步骤是：对式~\eqref{eq:2.13} 两边再取 ${\rm Hom}(-,U(1))$，左端得 ${\rm Hom}({}^L\neg\TT,U(1))$，右端由脚注的自反性得 $\Lambda_{\rm cochar}$；再用 $\Lambda_{\rm cochar}={\rm Hom}(U(1),\TT)$ 即得式~\eqref{eq:2.14}。物理上这正是 GNO 对偶：左端分类 $G$ 的磁荷，右端分类 $^L\neg G$ 的电荷。这一步用到的关键事实是 $\Lambda_{\rm cochar}$ 与 $\Lambda_{\rm char}$ 互为对偶格（见上文 ${}^L\neg\TT={}^L{\mathfrak t}/\Lambda_{\rm char}$ 的定义），因此它不是同义反复。
\end{yizhu}
